\documentclass[11pt]{article}
\usepackage[margin=1in]{geometry}
\usepackage[T1]{fontenc}
\usepackage[utf8]{inputenc}
\usepackage{graphicx}
\usepackage{booktabs}
\usepackage{tabularx}
\usepackage{array}
\usepackage{caption}
\usepackage{float}
\usepackage{amsmath}
\usepackage[hidelinks]{hyperref}
\usepackage{newtxtext,newtxmath}
\usepackage{setspace}
\setstretch{1.08}
\setlength{\parindent}{0pt}
\setlength{\parskip}{0.55em}
\setcounter{secnumdepth}{2}
\setcounter{topnumber}{3}
\setcounter{bottomnumber}{2}
\setcounter{totalnumber}{5}
\renewcommand{\topfraction}{0.90}
\renewcommand{\bottomfraction}{0.80}
\renewcommand{\textfraction}{0.10}
\renewcommand{\floatpagefraction}{0.85}
\captionsetup[figure]{labelformat=empty}
\title{Microplastic Identification Using AI-Driven Image Segmentation with Synthetic Ecological Context}
\author{Alex Dils\textsuperscript{1}, David Raymond\textsuperscript{2}, Jack Spottiswood\textsuperscript{1}, Samay Kodige\textsuperscript{1}\\
Dylan Karmin\textsuperscript{3}, Rikhil Kokal\textsuperscript{1}, Win Cowger\textsuperscript{4,*}, Christoph Sadée\textsuperscript{5,*}\\[0.35em]
\small \textsuperscript{1}\,University of California, Berkeley\\
\small \textsuperscript{2}\,Dartmouth College\\
\small \textsuperscript{3}\,Massachusetts Institute of Technology\\
\small \textsuperscript{4}\,Moore Institute for Plastic Pollution Research\\
\small \textsuperscript{5}\,Division of Computational Medicine, Stanford University\\
\small \textsuperscript{*}\,Co-senior authors}
\date{}
\begin{document}
\maketitle
\begin{abstract}
Conventional methods for microplastic identification in water samples are costly, slow, and dependent on specialized expertise. We present a deep learning segmentation framework for identifying microplastic foreground in microscopy images and demonstrate that synthetic ecological context can improve segmentation model performance. The workflow combines manually labeled laboratory microplastic images with generated inpainting examples selected for visible foreground change, non-empty masks, plausible object scale, and background diversity, and it contributes a curated image collection with manually segmented microplastic masks for supervised training and held-out ecological evaluation. In our results, adding verified synthetic examples improves segmentation across architectures by increasing exposure to diverse ecological scenarios while preserving pixel-level labels. The strongest synthetic-assisted model reaches Dice 0.817, IoU 0.704, and boundary F1 0.858 on the held-out ecological evaluation set, compared with Dice 0.743, IoU 0.619, and boundary F1 0.809 for the strongest real-only model. These results extend prior microplastic segmentation work by testing synthetic ecological context under matched held-out evaluation and show that high-quality synthetic examples can improve segmentation.
\end{abstract}

\section{Introduction}
Microplastics are a widespread ecological concern, with detection and monitoring needs shaped by large-scale plastic production, environmental persistence, food-safety risks, and heterogeneous particle types \cite{1,6,7}. Microplastic monitoring needs methods that are accurate enough for scientific screening yet practical enough to apply at scale. Manual microscopy can flag candidate particles, but it is slow and subjective, and its reliability depends heavily on the analyst \cite{2,3,7}. Spectroscopic techniques such as FTIR and Raman provide far more specific confirmation, but they are costly and low-throughput, which limits how widely they can be deployed \cite{2,3,7}. A useful image-based system should sit between these extremes: it should help researchers prioritize candidate particles, quantify their morphology, and reduce the volume of manual review required before confirmatory analysis.

We approach this problem with semantic segmentation, which has been used in biomedical image tasks such as liver-tumor segmentation and is commonly implemented with U-Net-style architectures \cite{4,18}. In this case, segmentation labels microplastic at the pixel level, which makes it possible to recover the morphological quantities that matter for microplastic characterization: area, perimeter, width, shape, and contact with surrounding debris. Because a mask directly estimates which pixels belong to a particle, it is naturally suited to morphology and burden estimation. Detection remains valuable for triage and counting, but when the goal is quantitative measurement and morphology-aware screening, segmentation is the more appropriate first task.

The central obstacle to training such a model is the scarcity of labeled real-world microscopy images, particularly across diverse imaging conditions. Samples are particularly limited in the ecological setting, even though that is where segmentation is hardest and most useful. Recent reviews also point to data standardization and sharing as recurring constraints for microplastics research \cite{5,8}. Synthetic inpainting offers a practical way to bridge this gap: it can place known particle masks into realistic ecological backgrounds, expanding the visual diversity of the training data while preserving ground-truth labels.

Recent microplastic segmentation studies demonstrate that the task is feasible: Park et al.\ developed MP-Net, a U-Net-derived model for fluorescence microscopy images of microplastics isolated from clams, and reported a mean F1-score of 0.736 and mean IoU of 0.617 \cite{28}. Yao et al.\ reported 0.690 mIoU for a lightweight multi-class MNv4-Conv-M-fpn model \cite{24}, and Xu and Wang reported UNet and UNet2plus mIoUs of 91.45\% and 91.08\% on an urban-water microplastic dataset \cite{25}. Those studies establish strong task-specific segmentation performance; the question here is whether synthetic ecological context improves transfer to held-out ecological microscopy under matched real-only and synthetic-assisted training conditions.

This motivates the central question of the study: does adding synthetic ecological context improve segmentation performance on real samples? To answer it rigorously, we evaluate the effect across several segmentation architectures evaluating on a held-out ecologically diverse test set, so that any observed benefit reflects a general property of the synthetic data rather than an artifact of a single model.

\section{Methods}
\subsection{Datasets}
To train and test the models we gathered and organized data into three image cohorts, each serving a distinct role in training or evaluation. Together, the assembled image collection contains 1,198 microscopy images, including 465 manually segmented image-mask pairs and 733 unlabeled ecological background images.

\begin{table}[htbp]
\centering
\begingroup
\setlength{\tabcolsep}{3pt}
\footnotesize
\begin{tabularx}{\textwidth}{@{}>{\raggedright\arraybackslash}p{0.13\textwidth}>{\raggedright\arraybackslash}p{0.22\textwidth}>{\raggedright\arraybackslash}p{0.08\textwidth}>{\raggedright\arraybackslash}p{0.08\textwidth}>{\raggedright\arraybackslash}X@{}}
\toprule
\textbf{Cohort} & \textbf{Source type} & \textbf{Images} & \textbf{Masks} & \textbf{Role in study} \\
\midrule
Cohort 1 & Laboratory microplastic microscopy & 368 & 368 & Real supervised training source and microplastic mask source for inpainting \\
Cohort 2 & Unlicensed ecological microscopy without microplastic & 733 & 0 & Background image source for synthetic ecological context generation \\
Cohort 3 (test set) & Unlicensed ecological microplastic microscopy & 97 & 97 & Held-out evaluation set for final segmentation metrics \\
\bottomrule
\end{tabularx}
\caption{Dataset composition and role of each cohort.}
\endgroup
\end{table}

Cohort 1 consists of microscopy images of laboratory-prepared microplastic samples derived from the Moore Institute for Plastic Pollution Research Microplastic Image Explorer, an OpenAnalysis/One4All resource for querying microplastic microscopy images by citation, color, morphology, and polymer type \cite{8,29}. In this study, each selected microplastic particle was manually segmented and converted into a binary mask, yielding a one-to-one correspondence between images and labels. This cohort serves two purposes: it supplies the labeled examples for real supervised training, and its masks provide the particle templates inserted into ecological backgrounds during inpainting.

Cohort 2 comprises unlicensed ecological microscopy images that were webscraped and verified to not contain microplastic. The images were used as unlabeled background images without annotated microplastics. Because these images capture realistic environmental clutter---sediment, bubbles, fibers, and organic matter---they serve as the background source into which known particle masks are inpainted, expanding the visual diversity of the training data while preserving ground-truth labels.

Cohort 3 is the held-out evaluation set and is never seen during training. It contains unlicensed ecological microplastic microscopy images that were verified to contain microplastic; each particle was manually segmented into a binary mask. Reserving this cohort for final evaluation ensures that all reported metrics reflect performance on real ecological samples rather than on synthetic data. The manual segmentations in Cohorts 1 and 3 convert the source microscopy images into reusable pixel-level labels for training, benchmarking, and morphology-aware evaluation.

\subsection{Synthetic Model}
The custom inpainting generator was initialized from the pretrained Stability AI Stable Diffusion 2 inpainting checkpoint \cite{30} and fine-tuned on Cohort 1 laboratory microplastic images with their dilated binary masks; the full checkpoint identifier is reported in the model-settings appendix. Each training item consisted of the microscopy image, its corresponding inpainting mask, and the masked-image conditioning input used by the Stable Diffusion inpainting formulation. Images and masks were resized to 512$\times$512 pixels, and the model was trained to reconstruct the masked microplastic foreground while using the unmasked image region as visual context. The held-out Cohort 3 ecological evaluation images were not used in any inpainting-model training step.

The fine-tuned component carried forward for generation was the Diffusers inpainting U-Net checkpoint. Training used seed 42, 100 epochs, batch size 1, gradient accumulation of 4, learning rate $1\times10^{-5}$, a constant learning-rate schedule with no warmup, AdamW optimization with $\beta_1=0.9$, $\beta_2=0.999$, weight decay 0.01, epsilon $10^{-8}$, maximum gradient norm 1.0, and fp16 mixed precision. Memory-efficient attention was enabled when available. Checkpoints were written every 500 steps and at the end of each epoch.

Before segmentation training, synthetic image/mask pairs were produced using generative inpainting \cite{13,14,15,16,17}. Real Cohort 1 microplastic masks were geometrically transformed and inserted into unlabeled Cohort 2 ecological microscopy backgrounds to create masked ecological composites. The geometric transforms included random rotation, horizontal and vertical flipping, isotropic scaling, mild anisotropic scaling, translation to a random valid background location, and shear. Each transformed mask was placed only where it remained inside the image frame, and the corresponding real microplastic crop was inserted into the ecological background to create a mask-guided editing composite. This mask-guided editing composite provided the spatial location and approximate particle structure for inpainting, while the binary transformed mask defined the region to be regenerated.

Stable Diffusion inpainting \cite{13,14,15} was then used to synthesize plausible microplastic foreground appearance inside the masked region while preserving the surrounding ecological background. For each candidate example, the input to the generator consisted of the ecological background, the mask-guided editing composite, the transformed binary mask, and a text prompt specifying a realistic microscope water sample with a visible microplastic fiber or fragment in the masked region. Negative prompts penalized cartoon-like output, text, watermarks, blur, unchanged backgrounds, and empty masks. The generation settings used 40--45 denoising steps, guidance scale 7.0--8.0, inpainting strength 0.99, full-mask diffusion, and 4-pixel mask dilation. The retained transformed mask served as the segmentation label for each generated image. Stable Diffusion inpainting was used to generate 10,000 candidate synthetic image/mask pairs. These generator outputs were not evaluated against the held-out Cohort 3 labels; they were used only to create candidate training pairs for the downstream segmentation experiment.

During Stable Diffusion inpainting, generation was restricted to pixels inside the supplied binary mask, while all pixels outside the mask were preserved from the original ecological background, ensuring that only the masked foreground region was modified.

Synthetic inpainting examples were quality-controlled by scoring each generated image with four mask-region change metrics: masked mean absolute difference versus the original ecological background, changed-pixel fraction versus the original background, masked mean absolute difference versus the mask-guided editing composite, and changed-pixel fraction versus the mask-guided editing composite. Foreground mask fraction was also checked to exclude empty, tiny, or implausibly large masks. Candidate images were ranked by a composite normalized QC score, and the final top-inpainting subset was selected with diversity caps of one image per ecological background and two images per source microplastic mask. This retained generated examples that visibly changed the masked region, preserved plausible microplastic scale, and avoided overrepresenting repeated backgrounds or source masks.
\begin{equation}
\label{eq:qc-score}
\begin{aligned}
Q_i
&= \sum_{k \in \mathcal{K}} z_{ik}
   - 0.25\,
   \frac{\left|m_i - \operatorname{median}(m)\right|}
        {q_{0.95,m} - q_{0.05,m}}, \\
z_{ik}
&= \frac{\operatorname{clip}\left(x_{ik}, q_{0.01,k}, q_{0.99,k}\right)
        - q_{0.01,k}}
       {q_{0.99,k} - q_{0.01,k}}, \\
\mathcal{K}
&= \{\text{masked MAD vs background},
      \text{changed fraction vs background}, \\
&\qquad \text{masked MAD vs composite},
      \text{changed fraction vs composite}\}, \\
m_i
&= \text{mask foreground fraction}.
\end{aligned}
\end{equation}

\subsection{Segmentation Model}
The retained benchmark evaluates U-Net \cite{18}, U-Net++ \cite{19}, DeepLabV3+ \cite{20}, FPN \cite{21}, and SegFormer-B2 \cite{22}.

Models are trained with Dice loss \cite{26} and binary cross-entropy. A fixed threshold of 0.5 means that pixels with predicted foreground probability at or above 0.5 are counted as microplastic foreground during evaluation. All models use deterministic seeds 13, 37, and 101; images are resized to 512$\times$512 during training; batch size is 8; the maximum training horizon is 80 epochs; the learning rate is $1\times10^{-4}$ with weight decay $1\times10^{-5}$; early stopping patience is 15 epochs; and mixed precision is enabled when CUDA is available. Training augmentations include horizontal and vertical flips, translation, crop, shear, and color jitter. The primary metrics are Dice and IoU; secondary metrics include boundary F1 \cite{27}, precision, recall, and area error. Validation scores are used only for early stopping and diagnostics. The central evidence comes from the held-out ecological evaluation set.
\subsection{Experimental Design}

\begin{figure}[H]
	\centering
	\includegraphics[width=0.82\textwidth,height=0.45\textheight,keepaspectratio]{figures/figure_3_workflow_pipeline.png}
	\caption{Figure 1. Synthetic ecological context generation and evaluation pipeline. A generator fine-tuned from Cohort 1 image-mask pairs is applied to Cohort 2 backgrounds, synthetic image-mask pairs train the segmenter, and performance is evaluated on held-out Cohort 3 images.}
\end{figure}

\begin{table}[htbp]
	\centering
	\small
	\begin{tabularx}{\textwidth}{>{\raggedright\arraybackslash}p{0.24\linewidth}>{\raggedright\arraybackslash}p{0.30\linewidth}>{\raggedright\arraybackslash}p{0.38\linewidth}}
		\toprule
		\textbf{Condition} & \textbf{Training data} & \textbf{Purpose} \\
		\midrule
		Real only & Labeled Cohort 1 images only & Baseline for supervised segmentation without synthetic ecological context \\
		Real + unfiltered inpainting & Real images plus all generated inpainting examples & Tests whether synthetic scale alone is sufficient \\
		Real + verified inpainting & Real images plus QC-filtered inpainting examples & Tests whether high-quality synthetic ecological context improves transfer \\
		Half real + top inpainting & Equal real and top-ranked inpainting examples & Tests whether a smaller set of best synthetic samples can outperform larger noisier synthetic sets \\
		\bottomrule
	\end{tabularx}
	\caption{Experimental training conditions.}
\end{table}

The experiment is designed to isolate the effect of synthetic ecological context on segmentation performance. Every condition uses the same real labeled images, the same evaluation set, the same segmentation threshold, and the same training schedule. The comparison changes only whether synthetic inpainting examples are added to the real training data and, in the quality-controlled condition, which generated images are admitted.

Synthetic examples are generated by placing transformed real microplastic masks into ecological backgrounds and using inpainting to synthesize foreground appearance inside the masked region. The retained mask is used as the label. For the verified condition, the hard exclusions are non-empty masks and diversity caps of one image per ecological background and two images per source microplastic mask; the remaining candidates are ranked by the normalized QC score in Equation~\ref{eq:qc-score} rather than by a single manual cutoff. This prevents the model from learning from no-op generations, oversized artifacts, or repeated backgrounds.


\section{Results}
The verified synthetic examples are qualitatively plausible ecological microplastic training pairs. The retained examples show visible microplastic-like foregrounds embedded in ecological microscopy backgrounds rather than empty masks or unchanged background patches. The examples in Figure 2 preserve the intended transformed mask geometry while placing fibers and fragments into visually varied debris fields. Across the displayed cases, the synthetic foreground remains spatially aligned with the transformed mask, and the surrounding backgrounds contain heterogeneous debris, texture, and illumination that are absent from the cleaner source masks. This supports the premise that the synthetic images add ecological context without breaking the pixel-level supervision needed for segmentation.

\begin{figure}[H]
	\centering
	\includegraphics[width=0.56\textwidth,height=0.42\textheight,keepaspectratio]{figures/figure_2_inference_grid.png}
	\caption{Figure 2. Representative source microplastic examples with original masks, transformed masks, and inpainted ecological images.}
\end{figure}

\begin{table}[htbp]
	\centering
	\small
	\begin{tabularx}{\textwidth}{>{\raggedright\arraybackslash}p{0.24\linewidth}>{\raggedright\arraybackslash}p{0.08\linewidth}>{\raggedright\arraybackslash}p{0.16\linewidth}>{\raggedright\arraybackslash}p{0.16\linewidth}>{\raggedright\arraybackslash}p{0.15\linewidth}>{\raggedright\arraybackslash}p{0.15\linewidth}}
		\toprule
		\textbf{Training condition} & \textbf{Runs} & \textbf{Dice} & \textbf{IoU} & \textbf{Boundary F1} & \textbf{Best validation Dice} \\
		\midrule
		Real only & 15 & 0.641 $\pm$ 0.161 & 0.519 $\pm$ 0.138 & 0.703 & 0.758 \\
		Real + unfiltered inpainting & 15 & 0.586 $\pm$ 0.143 & 0.462 $\pm$ 0.132 & 0.642 & \textbf{0.928} \\
		Real + verified inpainting & 15 & 0.735 $\pm$ 0.061 & 0.608 $\pm$ 0.058 & 0.797 & 0.846 \\
		Half real + top inpainting & 15 & \textbf{0.754 $\pm$ 0.028} & \textbf{0.627 $\pm$ 0.031} & \textbf{0.814} & 0.812 \\
		\bottomrule
\end{tabularx}
\caption{Held-out ecological segmentation metrics by training condition. Values are mean $\pm$ standard deviation across 15 model-seed runs; bold entries mark the best value in each metric column.}
\end{table}

\begin{figure}[H]
\centering
\includegraphics[width=0.76\textwidth,height=0.27\textheight,keepaspectratio]{figures/figure_4_condition_metrics.png}
\caption{Figure 3. Held-out ecological segmentation metrics by training condition. Error bars show standard deviations where reported in the manuscript table.}
\end{figure}

Adding verified synthetic examples to real training data improves held-out ecological segmentation compared with real-only training. The strongest verified-inpainting checkpoint is U-Net++ at seed 37, with Dice 0.817, IoU 0.704, boundary F1 0.858, precision 0.802, and recall 0.846. The strongest real-only checkpoint reaches Dice 0.743, IoU 0.619, and boundary F1 0.809, so the top synthetic-assisted run improves Dice by 0.074 and IoU by 0.085 over the strongest real-only run. The mean condition-level comparison shows the same direction of effect: real-only training reaches mean Dice 0.641 and mean IoU 0.519, whereas real plus verified inpainting reaches mean Dice 0.735 and mean IoU 0.608. The synthetic effect is not limited to a single architecture: all five retained semantic architectures improve under the verified-inpainting condition, with absolute Dice gains ranging from +0.050 for U-Net and FPN to +0.102 for U-Net++.

Quality-controlled synthetic data outperform unfiltered synthetic data, showing that filtering is essential for the performance gain. Mean Dice falls from 0.641 for real-only training to 0.586 under real plus unfiltered inpainting, while verified inpainting raises mean Dice to 0.735 and mean IoU to 0.608. The unfiltered condition also produces the highest best validation Dice, 0.928, but this validation strength does not transfer to the held-out ecological test set, suggesting that unfiltered generations can make the training and validation distribution easier without improving ecological robustness. The half real plus top-inpainting condition further supports this pattern, reaching Dice 0.754 across 15 model-seed runs and matching or exceeding the full verified-inpainting mean on the primary overlap metrics. These results indicate that the useful signal comes from filtered, diverse, visibly changed synthetic examples rather than from adding synthetic data indiscriminately.



\begin{table}[htbp]
\centering
\small
\begin{tabularx}{\textwidth}{>{\raggedright\arraybackslash}p{0.16\linewidth}>{\raggedright\arraybackslash}p{0.15\linewidth}>{\raggedright\arraybackslash}p{0.18\linewidth}>{\raggedright\arraybackslash}p{0.14\linewidth}>{\raggedright\arraybackslash}p{0.29\linewidth}}
\toprule
\textbf{Model} & \textbf{Real-only Dice} & \textbf{Verified-inpainting Dice} & \textbf{Absolute gain} & \textbf{Interpretation} \\
\midrule
U-Net++ & 0.704 & 0.806 & +0.102 & Largest gain; multiscale decoder benefits from added ecological variation \\
SegFormer-B2 & 0.728 & 0.789 & +0.061 & Transformer features improve with broader background diversity \\
DeepLabV3+ & 0.694 & 0.758 & +0.064 & Atrous context helps exploit synthetic foreground-background variation \\
FPN & 0.685 & 0.735 & +0.050 & Moderate gain from multiscale synthetic examples \\
U-Net & 0.660 & 0.710 & +0.050 & Improves but remains less robust than U-Net++ \\
\bottomrule
\end{tabularx}
\caption{Architecture-level Dice change after verified inpainting for the retained segmentation models.}
\end{table}

\iffalse
\begin{figure}[!htbp]
\centering
\includegraphics[width=0.80\textwidth,height=0.40\textheight,keepaspectratio]{figures/figure_6_best_checkpoint_comparison.png}
\caption{Figure 4. Strongest real-only checkpoint versus strongest synthetic-assisted checkpoint on held-out ecological evaluation metrics.}
\end{figure}
\fi

\section{Discussion}
The synthetic examples appear to help for three reasons. First, they expose the model to ecological backgrounds that are absent from clean laboratory images. This reduces false positives on debris because the model sees more non-microplastic clutter during training. Second, they preserve pixel-level masks, so the model can learn microplastic shape while seeing new background contexts. Third, quality filtering removes failed generations before they can corrupt the training distribution. These factors together convert synthetic data from a source of noise into a controlled form of domain expansion.

The improvement is greatest in boundary F1, which rises from 0.703 for real-only training to 0.797 for verified inpainting. This supports the interpretation that synthetic ecological context improves not only particle detection but also boundary localization. That matters for microplastic analysis because area, length, and morphology are derived from the mask boundary.

The half real + top inpainting condition is consistent with this interpretation: using equal counts of real and top-ranked synthetic examples reaches Dice 0.754, exceeding both the real-only mean and the unfiltered-inpainting mean.

As an external point of reference, MP-Net achieved mean F1 0.736 and mean IoU 0.617 on fluorescence microscopy images of microplastics isolated from clams \cite{28}. Our strongest synthetic-assisted checkpoint is numerically higher on Dice/F1 and IoU, but this is not a direct benchmark comparison because the imaging modality, sample source, label protocol, and test distribution differ.

These findings should be interpreted with several limitations. Generated images, although visually plausible, may not capture the full variability and complexity of real-world ecological scenarios. Synthetic foregrounds may also encode generator-specific artifacts that improve benchmark performance without improving field reliability.

The image collection may not encompass all environmental contexts where microplastics occur, which may limit generalization to new sites, imaging devices, particle morphologies, or preparation protocols. The held-out ecological test set is intentionally reserved for final evaluation, but it remains modest in size. Additional externally collected test sets would be needed to confirm robustness across laboratories, cameras, and sample preparation workflows.

\section{Conclusion}
This study demonstrates that quality-controlled synthetic ecological context can improve microplastic segmentation when labeled real-world data are limited. It also contributes a curated set of manually segmented microplastic microscopy images, providing paired image-mask labels for real supervised training and held-out ecological evaluation. Across the retained architectures, verified synthetic examples improved held-out ecological segmentation relative to real-only training, with the strongest synthetic-assisted model outperforming the strongest real-only model by 0.074 Dice and 0.085 IoU.

The consistent gains across the retained model families support the central conclusion that synthetic ecological context can complement real labeled images when candidate generations are filtered before training.

Future work should evaluate the approach across additional imaging platforms, environmental settings, and microplastic morphologies, as well as investigate more advanced generative methods and larger-scale ecological datasets. Integrating quality-controlled synthetic data with deep learning segmentation may help enable faster, more scalable, and more accessible microplastic monitoring workflows for environmental research and pollution assessment.
\section*{Data and Code Availability}
The project code and training pipeline are available in the \href{https://github.com/axel-slid/synthetic-microplastic-detection}{project GitHub repository}. Cohort 1 source microscopy imagery is available through the \href{https://www.openanalysis.org/microplastic_image_explorer/}{OpenAnalysis Microplastic Image Explorer} \cite{29}, and the Cohort 1 microplastic image-mask pairs are available through Harvard Dataverse at \href{https://doi.org/10.7910/DVN/WG3OTM}{doi:10.7910/DVN/WG3OTM} \cite{31}. Cohorts 2 and 3 were assembled from unlicensed microscopy images rather than from a formal dataset. The project image collection includes the manually segmented microplastic image-mask pairs used for real supervised training and held-out ecological evaluation, along with the ecological background images used for synthetic context generation. All publication-ready results should be regenerated from the locked evaluation files and registered checkpoints.
\section*{Acknowledgments}
We thank the Moore Institute for Plastic Pollution Research for providing microplastic imagery and supporting data resources used in this work.
\section*{Generative AI Statement}
Generative AI tools were used to assist with preparation of this manuscript, including drafting, editing, and language refinement. The authors reviewed, revised, and approved all content and take full responsibility for the final manuscript.
\begin{thebibliography}{31}
\bibitem{1} Ziani, K. et al. (2023) Microplastics: A real global threat for environment and food safety: A state of the art review. Nutrients 15(3):617. doi:10.3390/nu15030617.
\bibitem{2} Stock, F. et al. (2020) Pitfalls and limitations in microplastics analyses. Plastics in the Aquatic Environment -- Part I. The Handbook of Environmental Chemistry, vol. 111. Springer, Cham. doi:10.1007/698\_2020\_654.
\bibitem{3} Primpke, S., Christiansen, S.H., Cowger, W., et al. (2020) Critical assessment of analytical methods for the harmonized and cost-efficient analysis of microplastics. Applied Spectroscopy 74(9):1012--1047. doi:10.1177/0003702820921465.
\bibitem{4} Sabir, M.W. et al. (2022) Segmentation of liver tumor in CT scan using ResU-Net. Applied Sciences 12(17):8650. doi:10.3390/app12178650.
\bibitem{5} Zhang, Y., Zhang, D., and Zhang, Z. (2023) A critical review on artificial intelligence-based microplastics imaging technology: Recent advances, hot-spots, and challenges. International Journal of Environmental Research and Public Health 20:1150. doi:10.3390/ijerph20021150.
\bibitem{6} Geyer, R., Jambeck, J.R., and Law, K.L. (2017) Production, use, and fate of all plastics ever made. Science Advances 3:e1700782. doi:10.1126/sciadv.1700782.
\bibitem{7} Hidalgo-Ruz, V., Gutow, L., Thompson, R., and Thiel, M. (2012) Microplastics in the marine environment: A review of the methods used for identification and quantification. Environmental Science \& Technology 46:3060--3075. doi:10.1021/es2031505.
\bibitem{8} Sherrod, H. et al. (2024) One4All: An open source portal to validate and share microplastics data and beyond. Journal of Open Source Software 9(99):6715. doi:10.21105/joss.06715.
\bibitem{9} Brandt, J., Mattsson, K., and Hassellöv, M. (2021) Deep learning for reconstructing low-quality FTIR and Raman spectra: A case study in microplastic analyses. Analytical Chemistry 93:16360--16368. doi:10.1021/acs.analchem.1c02618.
\bibitem{10} Yurtsever, M. and Yurtsever, U. (2019) Use of a convolutional neural network for the classification of microbeads in urban wastewater. Chemosphere 216:271--280. doi:10.1016/j.chemosphere.2018.10.084.
\bibitem{11} Shorten, C. and Khoshgoftaar, T.M. (2019) A survey on image data augmentation for deep learning. Journal of Big Data 6:60. doi:10.1186/s40537-019-0197-0.
\bibitem{12} Isola, P., Zhu, J.-Y., Zhou, T., and Efros, A.A. (2017) Image-to-image translation with conditional adversarial networks. CVPR. doi:10.1109/CVPR.2017.632.
\bibitem{13} Rombach, R., Blattmann, A., Lorenz, D., Esser, P., and Ommer, B. (2022) High-resolution image synthesis with latent diffusion models. CVPR. doi:10.1109/CVPR52688.2022.01042.
\bibitem{14} Podell, D. et al. (2023) SDXL: Improving latent diffusion models for high-resolution image synthesis. arXiv:2307.01952.
\bibitem{15} Lugmayr, A. et al. (2022) RePaint: Inpainting using denoising diffusion probabilistic models. CVPR. doi:10.1109/CVPR52688.2022.01117.
\bibitem{16} Suvorov, R. et al. (2022) Resolution-robust large mask inpainting with Fourier convolutions. WACV. doi:10.1109/WACV51458.2022.00323.
\bibitem{17} Li, W., Lin, Z., Zhou, K., Qi, L., Wang, Y., and Jia, J. (2022) MAT: Mask-aware transformer for large hole image inpainting. CVPR. doi:10.1109/CVPR52688.2022.01049.
\bibitem{18} Ronneberger, O., Fischer, P., and Brox, T. (2015) U-Net: Convolutional networks for biomedical image segmentation. MICCAI, 234--241. doi:10.1007/978-3-319-24574-4\_28.
\bibitem{19} Zhou, Z., Siddiquee, M.M.R., Tajbakhsh, N., and Liang, J. (2018) UNet++: A Nested U-Net architecture for medical image segmentation. Deep Learning in Medical Image Analysis and Multimodal Learning for Clinical Decision Support, 3--11. doi:10.1007/978-3-030-00889-5\_1.
\bibitem{20} Chen, L.-C., Zhu, Y., Papandreou, G., Schroff, F., and Adam, H. (2018) Encoder-decoder with atrous separable convolution for semantic image segmentation. ECCV, 833--851. doi:10.1007/978-3-030-01234-2\_49.
\bibitem{21} Lin, T.-Y., Dollár, P., Girshick, R., He, K., Hariharan, B., and Belongie, S. (2017) Feature pyramid networks for object detection. CVPR, 2117--2125. doi:10.1109/CVPR.2017.106.
\bibitem{22} Xie, E., Wang, W., Yu, Z., Anandkumar, A., Alvarez, J.M., and Luo, P. (2021) SegFormer: Simple and efficient design for semantic segmentation with transformers. NeurIPS. arXiv:2105.15203.
\bibitem{23} Jocher, G. and Qiu, J. (2024) Ultralytics YOLO11. Version 11.0.0. https://github.com/ultralytics/ultralytics
\bibitem{24} Yao, Y., Xu, W., and Fan, H. (2025) A Deep Learning Approach for Microplastic Segmentation in Microscopic Images. Toxics 13(12):1018. doi:10.3390/toxics13121018.
\bibitem{25} Xu, J. and Wang, Z. (2024) Efficient and accurate microplastics identification and segmentation in urban waters using convolutional neural networks. Science of the Total Environment 911:168696. doi:10.1016/j.scitotenv.2023.168696.
\bibitem{26} Milletari, F., Navab, N., and Ahmadi, S.-A. (2016) V-Net: Fully convolutional neural networks for volumetric medical image segmentation. 3DV, 565--571. doi:10.1109/3DV.2016.79.
\bibitem{27} Martin, D.R., Fowlkes, C.C., and Malik, J. (2004) Learning to detect natural image boundaries using local brightness, color, and texture cues. IEEE Transactions on Pattern Analysis and Machine Intelligence 26(5):530--549. doi:10.1109/TPAMI.2004.1273918.
\bibitem{28} Park, H.-m., Park, S., de Guzman, M.K., Baek, J.Y., Cirkovic Velickovic, T., Van Messem, A., and De Neve, W. (2022) MP-Net: Deep learning-based segmentation for fluorescence microscopy images of microplastics isolated from clams. PLOS ONE 17(6):e0269449. doi:10.1371/journal.pone.0269449.
\bibitem{29} Sherrod, H. and Cowger, W. (2024) Microplastic Image Explorer Tutorial. One4All 0.5. Moore Institute for Plastic Pollution Research / OpenAnalysis. https://www.openanalysis.org/One4All/articles/imageexplorer.html; application: https://www.openanalysis.org/microplastic\_image\_explorer/.
\bibitem{30} Lopez, J. (2022) Stable Diffusion 2.0 Release. Stability AI. https://stability.ai/news-updates/stable-diffusion-v2-release.
\bibitem{31} Anonymous. (2024) Cohort 1: Microplastic Images and Microplastic Masks (368 images). Harvard Dataverse. doi:10.7910/DVN/WG3OTM.
\end{thebibliography}
\clearpage
\appendix
\renewcommand{\thesection}{Appendix \Alph{section}}
\renewcommand{\thesubsection}{\Alph{section}.\arabic{subsection}}
\section{Model Settings}
\subsection{Shared Study Settings}
All images were resized to 512$\times$512 pixels. The random seeds were 13, 37, and 101. The primary evaluation split was the locked C3-clean ecological test set, and the primary threshold for binary masks was 0.5. Semantic segmentation models were trained with Dice loss plus binary cross entropy using the AdamW optimizer. The batch size was 8, the maximum training horizon was 80 epochs, the learning rate was $1\times10^{-4}$, and the weight decay was $1\times10^{-5}$. Early stopping patience was 15 epochs. AMP was enabled when CUDA was available. Synthetic examples were split into 80\% training and 20\% validation subsets.

\subsection{Stable Diffusion Inpainting Model}
The Stable Diffusion inpainting model was used for synthetic ecological microplastic image generation. The generator family was Stable Diffusion inpainting. The primary model ID was diffusers / stable-diffusion-xl-1.0-inpainting-0.1, and the Stable Diffusion 2 inpainting model ID was stabilityai / stable-diffusion-2-inpainting. The image size was 512$\times$512 pixels.

For SDXL inpainting, the positive prompt was: ``a realistic microscope ecological water sample containing one visible translucent colored microplastic fiber or fragment inside the masked region, natural debris, sharp focus.'' The negative prompt was: ``cartoon, illustration, fake texture, text, watermark, blurry, unchanged background, empty mask, pipe, thick tube.'' The SDXL guidance scale was 7.0, the SDXL inference step count was 40, and the SDXL strength was 0.99. The SD2 guidance scale was 7.5, the SD2 inference step count was 45, and the SD2 strength was 0.99.

Seeded object initialization was enabled with a seeded object mix of 0.65. Diffusion blend was enabled, the diffusion mask mode was full, and inpaint masks were dilated by 4 px. Low-change generations were rejected. The minimum masked MAD was 12.0, the minimum changed-pixel fraction was 0.35, the minimum background masked MAD was 20.0, and the minimum background changed-pixel fraction was 0.50. The maximum generation attempt multiplier was 12. The insertion area fraction ranged from 0.006 to 0.05. Object opacity was 0.95, alpha feather radius was 0.7, and color jitter used brightness from 0.85 to 1.15 and contrast from 0.90 to 1.15.

\subsection{Semantic Segmentation Models}
\paragraph{U-Net ResNet34.}
The U-Net ResNet34 model was named smp\_unet\_resnet34. It was a semantic segmentation model implemented with segmentation-models-pytorch, using the Unet architecture with a resnet34 encoder. The model used 3 input channels and produced 1 binary foreground mask class. Training used logits, and evaluation applied a sigmoid activation followed by the 0.5 threshold.

\paragraph{U-Net++ EfficientNet-B4.}
The U-Net++ EfficientNet-B4 model was named smp\_unetpp\_effb4. It was a semantic segmentation model implemented with segmentation-models-pytorch, using the UnetPlusPlus architecture with an efficientnet-b4 encoder. The model used 3 input channels and produced 1 binary foreground mask class. Training used logits, and evaluation applied a sigmoid activation followed by the 0.5 threshold.

\paragraph{DeepLabV3+ ResNet50.}
The DeepLabV3+ ResNet50 model was named smp\_deeplabv3plus\_resnet50. It was a semantic segmentation model implemented with segmentation-models-pytorch, using the DeepLabV3Plus architecture with a resnet50 encoder. The model used 3 input channels and produced 1 binary foreground mask class. Training used logits, and evaluation applied a sigmoid activation followed by the 0.5 threshold.

\paragraph{FPN EfficientNet-B3.}
The FPN EfficientNet-B3 model was named smp\_fpn\_effb3. It was a semantic segmentation model implemented with segmentation-models-pytorch, using the FPN architecture with an efficientnet-b3 encoder. The model used 3 input channels and produced 1 binary foreground mask class. Training used logits, and evaluation applied a sigmoid activation followed by the 0.5 threshold.

\paragraph{SegFormer-B2.}
The SegFormer-B2 model was named segformer\_b2. It was a semantic segmentation model implemented with Hugging Face transformers, using the nvidia/segformer-b2-finetuned-ade-512-512 model ID. The model used 3 input channels. The pretrained ADE head was replaced with a binary segmentation head that produced 1 binary foreground mask class. The output logits were upsampled to the 512$\times$512 mask size. Training used logits, and evaluation applied a sigmoid activation followed by the 0.5 threshold.
\end{document}
